\chapter*{Terminological Index}
\phantomsection
\addcontentsline{toc}{chapter}{Terminological Index}

The reference numbers indicate successively the paragraph and the number (or, exceptionally, the exercise).

Adjoint (right, left) of a homomorphism: 1, 8.

Adjoint of a semilinear map: 7, 3.

Clifford algebra of a quadratic form: 9, 1.

Alternating (matrix): 3, 1.

Angle of two half-lines: 10, 3, and exercise 6.

Angle between two rays in an affine plane: 10, 3.

Angle between two lines: 10, 3 and exercise 2.

Angle of two lines in an affine plane: 10, 3.

Angle between two directed lines: 10, exercise 2.

Angle between two vectors in a plane: 10, 3.

Angle between two vectors in a space of dimension \(\geqslant 2:10,3\) .

Right angle: 10, 3 and exercise 2.

Angle of a rotation: 10, 3.

Angle of a direct similarity: 10, exercise 6.

Straight angle: 10, 3 and exercise 6.

Ring of types of quadratic forms: 8, 3.

Witt ring: 8, 3.

Antiautomorphism of a ring: 1, 2.

Principal antiautomorphism of a Clifford algebra: 9, 1.

Antihermitian (form): 3, 1.

Antihermitian (matrix): 3, 1.

Antisymmetric (matrix): 3, 1.

Bilinear map: 1, 1.

Degenerate bilinear map (on the right, on the left): 1, 1.

Non-degenerate bilinear map: 1, 1.

Nondegenerate bilinear map associated with a bilinear map: 1, 3.

Bilinear map obtained by extension of scalars from a bilinear map: 1, 4.

Canonical map of a module into its Clifford algebra: 9, 1.

Canonical mapping from the set of Hermitian endomorphisms to the set of Hermitian forms: 7, 3.

Canonical mapping of the set \(\mathfrak{A}\) on \(\mathrm{S}^{+}/\mathrm{H}:10,3\) .

Canonical mapping of the set \(\mathfrak{A}\) on \(O^{+}:10,3\) .

Right (left) linear map associated with a bilinear map: 1, 1.

Right (left) semilinear map associated with a sesquilinear form: 1, 6.

Right (left) sesquilinear map for an antiautomorphism: 1, 2. Degenerate (right, left) sesquilinear map: 1, 2. Nondegenerate sesquilinear map: 1, 2. Nondegenerate sesquilinear map associated with a sesquilinear map: 1, 3. Sesquilinear map obtained by extension of scalars from a sesquilinear map: 1, 4. Associated (linear map) with a bilinear map: 1, 1. Associated (semilinear map) with a sesquilinear form: 1, 6. Associated (bilinear form) with a quadratic form: 3, 4. Orthogonal automorphism: 6, 2. Orthogonal automorphism in n variables: 6, 2. Principal automorphism of a Clifford algebra: 9, 1. Symplectic automorphism: 5, 3. Symplectic automorphism in 2m variables: 5, 3. Unitary automorphism: 6, 2. Unitary automorphism in n variables: 6, 2. Axis of an affine linear complex: 10, exerc. 16. Orthogonal basis: 6, 1. Orthonormal basis: 6, 1. Symplectic basis: 5, 1. Bilinear (map): 1, 1. Bilinear (form): 1, 6. Bimodule: 1, 1. Canonical: see Canonical map and Canonical homomorphism. Center of a quadric: 6, exerc. 25 and 10, exerc. 12. Unit circle: 10, exerc. 2. Chasles (relation of): 10, 3. Clifford (algebra of): 9, 1. Clifford (group of): 9, 5. Linear complex: 10, exerc. 16. Condition (C): 6, 1. Condition (C'): 6, 1. Condition (T): 4, 2. Isotropic cone: 6, exerc. 23. Conformal (group): 10, exerc. 14. Affine conic: 6, exerc. 25. Projective conic: 6, exerc. 23. Conjugate (affine linear varieties): 6, exerc. 25. Conjugate (projective linear varieties): 6, exerc. 23. Cosine of an angle: 10, 3. Cotangent of an angle: 10, 3. Witt decomposition: 4, 2. Degenerate (bilinear map) on the right (left): 1, 1. Degenerate (sesquilinear map) on the right (left): 1, 2. Degenerate (affine conic): 6, exerc. 25. Degenerate (projective conic): 6, exerc. 23. Degenerate (affine quadric): 6, exerc. 25. Degenerate (projective quadric): 6, exerc. 23. Half-line: 10, 3. Closed half-line: 10, 3. Isotropic half-line: 10, 3. Open half-line: 10, 3. Displacement: 6, 6. Dimension of a quadratic form: 8, 1. Direct (similarity): 6, 5. Discriminant of a sesquilinear form with respect to a system of elements: 2. Distance: 7, 1. Right (angle): 10, 3. Pointed line: 10, exerc. 2.

Odd element of a Clifford algebra: 9, 1. Isotropic element: 4, 1. Even element of a Clifford algebra: 9, 1. Singular element: 4, 1. Elements orthogonal with respect to a sesquilinear map: 1, 3. Elements orthogonal with respect to a quadratic form: 3, 4. Hermitian endomorphism: 7, 3. Normal endomorphism: 7, 3. ε-Hermitian (form): 3, 1. ε-Hermitian (matrix): 3, 1. Equivalent (quadratic forms): 3, 4. Equivalent (sesquilinear forms): 1, 6. Definition space of a quadratic form: 8, 1. Euclidean space: 6, 6. Hermitian space: 6, 6. Oriented space: 10, 3. Euclidean (space): 6, 6. Extension of a sesquilinear form to a tensor power: 1, 9. Extension of a sesquilinear form to an exterior power: 1, 9. External (direct sum) of quadratic forms: 3, 4. External (direct sum) of quadratic modules: 3, 4. Weakly orthogonal (subspaces): 3, exerc. 11. Closed (angular sector): 10, 4. Closed (half-line): 10, 3. Cosine function: 10, 3. Cotangent function: 10, 3. Sine function: 10, 3. Tangent function: 10, 3. Trigonometric function: 10, 3. Anti-Hermitian form: 3, 1. Bilinear form: 1, 6. Bilinear form associated with a quadratic form: 3, 4. Hermitian bilinear form: 3, 1. Hermitian form: 3, 1. Negative (positive) Hermitian form: 7, 1. Inverse form of a bilinear (sesquilinear) form: 1, 7. Metric form: 6, 6. Neutral form: 4, 2 and 8, 1. Quadratic form: 3, 4. Negative (positive) quadratic form: 7, 1. Non-degenerate quadratic form: 3, 4. Quadratic form obtained by extension of scalars from a quadratic form: 3, 4. Sesquilinear form: 1, 6. Equivalent quadratic forms: 3, 4. Equivalent sesquilinear forms: 1, 6. Gram-Schmidt (orthogonalization process of): 6, 1. Clifford group: 9, 5. Reduced Clifford group: 9, 5. Special Clifford group: 9, 5. Rotation group: 9, 5. Group of types of quadratic forms: 8, 2. Witt group: 8, 2. Orthogonal group associated with Q: 6, 2. Orthogonal group in n variables: 6, 2. Special orthogonal group: 6, 2. Special unitary group: 6, 2. Symplectic group associated with Φ: 5, 3. Symplectic group in 2m variables: 5, 3. Unitary group associated with Φ: 6, 2.

\begin{Shaded}
\begin{Highlighting}[]
\NormalTok{Groupe unitaire à n variables : 6, 2.}
\NormalTok{Hermitien (endomorphisme) : 7, 3.}
\NormalTok{Hermitien (espace) : 6, 6.}
\NormalTok{Hermitienne (forme) : 3, 1.}
\NormalTok{Hermitienne (matrice) : 3, 1.}
\NormalTok{Homomorphisme canonique de l\textquotesingle{}algèbre de Clifford d\textquotesingle{}un sous{-}module de E dans l\textquotesingle{}algèbre de Clifford de E : 9, 1.}
\NormalTok{Homomorphisme métrique : 4, 3.}
\NormalTok{Hyperplan radical de deux sphères : 10, exerc. 12.}
\NormalTok{Image réciproque d\textquotesingle{}une application bilinéaire : 1, 1.}
\NormalTok{Image réciproque d\textquotesingle{}une application sesquilinéaire : 1, 2.}
\NormalTok{Impair (élément) dans une algèbre de Clifford : 9, 1.}
\NormalTok{Indice d\textquotesingle{}une forme quadratique (sesquilinéaire) : 4, 2.}
\NormalTok{Invariant de Dickson : 9, exerc. 9.}
\NormalTok{Inverse (forme) : 1, 7.}
\NormalTok{Inverse (similitude) : 6, 5.}
\NormalTok{Inversion : 10, exerc. 13.}
\NormalTok{Inversion de sphère C : 10, exerc. 13.}
\NormalTok{Involution dans un groupe linéaire : 6, 3.}
\NormalTok{Isotrope (demi{-}droite) : 10, 3.}
\NormalTok{Isotrope (élément) : 4, 1.}
\NormalTok{Isotrope (sous{-}module) : 4, 1.}
\NormalTok{Isotrope (variété linéaire) : 6, 6.}
\NormalTok{Laguerre (formule de) : 10, exerc. 5.}
\NormalTok{Loi d\textquotesingle{}inertie : 7, 2.}
\NormalTok{Longueur d\textquotesingle{}un vecteur : 10, 3.}
\NormalTok{Matrice alternée : 3, 1.}
\NormalTok{Matrice antihermitienne : 3, 1.}
\NormalTok{Matrice antisymétrique : 3, 1.}
\NormalTok{Matrice d\textquotesingle{}une application satisfaisant à (G), (D), (GD) ou (DG) : 1, 10.}
\NormalTok{Matrice d\textquotesingle{}une application bilinéaire : 1, 1.}
\NormalTok{Matrice d\textquotesingle{}une application sesquilinéaire : 1, 2.}
\NormalTok{Matrice ε{-}hermitienne : 3, 1.}
\NormalTok{Matrice hermitienne : 3, 1.}
\NormalTok{Matrice normale : 7, exerc. 17.}
\NormalTok{Matrice orthogonale : 6, 2.}
\NormalTok{Matrice symétrique : 3, 1.}
\NormalTok{Matrice symplectique : 5, 3.}
\NormalTok{Matrice unitaire : 6, 2.}
\NormalTok{Métrique (forme) : 6, 6.}
\NormalTok{Métrique (homomorphisme) : 4, 3.}
\NormalTok{Module quadratique : 3, 4.}
\NormalTok{Multiplicateur d\textquotesingle{}une similitude : 6, 5 et 6.}
\NormalTok{Négatif (n{-}vecteur) : 10, 3.}
\NormalTok{Négative (forme hermitienne) : 7, 1.}
\NormalTok{Négative (forme quadratique) : 7, 1.}
\NormalTok{Neutre (forme quadratique) : 8, 1.}
\NormalTok{Neutre (forme sesquilinéaire) : 4, 2.}
\NormalTok{Non dégénérée (application bilinéaire) : 1, 1.}
\NormalTok{Non dégénérée (application sesquilinéaire) : 1, 2.}
\NormalTok{Non dégénérée (forme quadratique) : 3, 4.}
\NormalTok{Normal (endomorphisme) : 7, 3.}
\NormalTok{Norme spinorielle : 9, 5.}
\NormalTok{Noyau d\textquotesingle{}un module quadratique : 3, 4.}
\NormalTok{n{-}vecteur négatif (positif) : 10, 3.}
\NormalTok{Orientation d\textquotesingle{}un espace : 10, 3.}
\NormalTok{Orienté (espace) : 10, 3.}
\NormalTok{Orthogonal (automorphisme) : 6, 2.}
\NormalTok{Orthogonal (groupe) : 6, 2.}
\NormalTok{Orthogonal (groupe spécial) : 6, 2.}
\end{Highlighting}
\end{Shaded}

Orthogonal (projector): 6, 3. Orthogonal (submodule) to a submodule: 1, 3. Orthogonal (vector) to a linear manifold: 6, 6. Orthogonal (basis): 6, 1. Orthogonal (matrix): 6, 2. Orthogonal (projection) onto a linear manifold: 6, 6. Orthogonal (transformation): 6, 2. Orthogonal (manifold) to a linear manifold: 6, 6. Orthogonal (parts): 1, 3 and 3, 4. Orthogonal (linear manifolds): 6, 6. Orthogonalization (process of) of Gram-Schmidt: 6, 1. Orthogonal (elements): 1, 3 and 3, 4. Orthonormal (basis): 6, 1. Open (angular sector): 10, 4. Open (half-line): 10, 3. Even (element) in a Clifford algebra: 9, 1. Orthogonal parts: 1, 3 and 3, 4. Perpendicular (linear manifolds): 6, exerc. 22. Pfaffian of an alternating matrix: 5, 2. Flat (angle): 10, 3. Point of view of a stereographic projection: 10, exerc. 14. Polar of a linear manifold with respect to a quadric: 6, exerc. 23 and 25. Pole of a hyperplane with respect to a quadric: 6, exerc. 23 and 25. Pole of an inversion: 10, exerc. 13. Positive (n-vector): 10, 3. Positive (Hermitian form): 7, 1. Positive (quadratic form): 7, 1. Principal (antiautomorphism): 9, 1. Principal (automorphism): 9, 1. Tensor product of quadratic forms: 8, 3. Tensor product of sesquilinear forms: 1, 9. Orthogonal projector: 6, 3. Orthogonal projection onto a linear manifold: 6, 6. Stereographic projection: 10, exerc. 14. Pseudo-discriminant: 9, exerc. 9. Power of a point with respect to a sphere: 10, exerc. 12. Power of an inversion: 10, exerc. 13. Pythagoras (theorem of): 6, 6. Quadratic (form): 3, 4. Quadratic (module): 3, 4. Affine quadric: 6, exerc. 25. Projective quadric: 6, exerc. 23. Rank of a bilinear (sesquilinear) form: 1, 6. Radius of a sphere: 10, exerc. 12. Reduced (Clifford group): 9, 5. Chasles relation: 10, 3. Spinor representations: 9, 4. Rotation: 9, 5. Rotation of angle φ: 10, 3 and exerc. 2. Rotations (group of): 9, 5. Closed angular sector: 10, 4. Open angular sector: 10, 4. Flat angular sector: 10, exerc. 8. Reentrant angular sector: 10, exerc. 8. Salient angular sector: 10, exerc. 8. Semi-spinor: 9, 4. Sesquilinear (map): 1, 2. Sesquilinear (form): 1, 6. Signature of a Hermitian form: 7, 2. Similarity: 6, 5 and 6.

Direct similarity: 6, 5 and 9, exercise 9.

Inverse similarity: 6, 5.

Singular (element): 4, 1.

Singular (submodule): 4, 1.

Sine of an angle: 10, 3.

Direct sum of bilinear (sesquilinear) maps: 1, 3.

External direct sum of quadratic forms (quadratic modules): 3, 4.

Isotropic submodule: 4, 1.

Orthogonal submodule to a submodule: 1, 3.

Singular submodule: 4, 1.

Totally isotropic submodule: 4, 1.

Totally singular submodule: 4, 1.

Sphere: 10, exercise 12.

Orthogonal spheres: 10, exercise 12.

Spinor: 9, 4.

Direct sequence of half-lines: 10, 4.

Symmetry with respect to a hyperplane: 6, 4.

Symmetry with respect to a vector subspace: 6, 3.

Symmetry with respect to an affine linear manifold: 6, 6.

Symmetric (matrix): 3, 1.

Symplectic (automorphism): 5, 3.

Symplectic (base): 5, 1.

Symplectic (group): 5, 3.

Symplectic (matrix): 5, 3.

Symplectic (transformation): 5, 3.

Tangent of an angle: 10, 3.

Tangent (linear variety) to a quadric: 6, exercises 23 and 25.

Pythagorean theorem: 6, 6.

Witt's theorem: 4, 3.

Totally isotropic (submodule): 4, 1.

Totally orthogonal (submodule) to a submodule: 1, 3.

Totally orthogonal (variety) to a linear variety: 6, 6.

Totally singular (submodule): 4, 1.

Orthogonal transformation: 6, 2.

Symplectic transformation: 5, 3.

Unitary transformation: 6, 2.

Trigonometric (function): 10, 3.

Type of a quadratic form: 8, 1.

Unitary (automorphism): 6, 2.

Unitary (group): 6, 2.

Unitary (special group): 6, 2.

Unitary (matrix): 6, 2.

Unitary (transformation): 6, 2.

Isotropic linear manifold: 6, 6.

Linear manifold orthogonal to a linear manifold: 6, 6.

Totally isotropic linear manifold: 6, 6.

Orthogonal linear manifolds: 6, 6.

Vector orthogonal to a linear manifold: 6, 6.

Witt (ring of): 8, 3.

Witt (decomposition of): 4, 2.

Witt (group of): 8, 2.

Witt (theorem of): 4, 3.

